\section{The prime 2}\label{sec:prime-two}

This section proves the lower bound for $\vG2(T)$ that gives the term $c_{2,\mathcal C}\log2$
of $\aconst_\sigma(q)$ (\cref{thm:prime-two}); it is used in \cref{thm:arithmetic}. The prime $2$
needs its own argument because $v_2(y)=-1$: the poles $x_k$ are not $2$-adically integral. In the
coordinate $z=2u$ they become the odd integers $z_k=-2k-1$, while the summation points become the
even integers $2m$, and the gain comes from treating the two parities separately.

\subsection{The rescaled determinant and its \texorpdfstring{$2$}{2}-adic expansion}\label{subsec:rescaled}

Put $P_K(z)=\prod_{k<K}(z-z_k)$, $N^{(2)}(z)=\prod_{j=1}^A(z-2j)^2$ and $F(z)=zN^{(2)}(z)/P_K(z)$.
Since $D_K(u)=2^{-K}P_K(2u)$ and $N_A(u)=2^{-2A}N^{(2)}(2u)$,
\begin{equation}\label{eq:phi-F}
  \varphi(u)=2^{K-1-2A}F(2u).
\end{equation}
On the space spanned by $\Q[z]$ and the fractions $1/(z-z_k)$ define the affine functional
\[
  \nu_X(G)=\mom\bigl(m\mapsto G'(2m)\bigr)\ (G\in\Q[z]),\qquad
  \nu_X\Bigl(\frac1{z-z_k}\Bigr)=-\sigma^ky^{2k+1}(X-\tau_k).
\]
If $R(u)=G(2u)$ then $R'(m)=2G'(2m)$, and $1/(u-x_k)=2/(2u-z_k)$; comparing
with~\eqref{eq:functional} gives $\mu_X(R)=2\nu_X(G)$ on both kinds of generators. With
$H(X)=[\nu_X(Fz^{i+j})]_{0\le i,j<h}$ and~\eqref{eq:phi-F}, the $(i,j)$ entry of the matrix
defining $\Delta$ is $2^{K-2A-i-j}H_{ij}$, hence
\begin{equation}\label{eq:Delta-H}
  \Delta(X)=2^{(K-2A)h-h(h-1)}\det H(X).
\end{equation}

The series $\xi_2=\sum_{j\ge0}\sigma^jy^{-2j-1}/(2j+1)^2$ converges in $\Z_2$, because the $j$-th
term has valuation $2j+1$. Repeating the computation in the proof of \cref{lem:pole-series} in
$\Q_2$, where $\sum_mw^m(\cdots)$ converges because $v_2(w)=2$, gives
\[
  \sum_{m\ge1}w^m\frac{\mathrm d}{\mathrm dz}\Bigl(\frac1{z-z_k}\Bigr)\Big|_{z=2m}
  =-\sigma^ky^{2k+1}(\xi_2-\tau_k).
\]
Hence $\nu_X(1/(z-z_k))=\sum_mw^m(\cdots)-\sigma^ky^{2k+1}(X-\xi_2)$. Writing
$Fz^s=Q_s+\sum_kr_kz_k^s/(z-z_k)$ with $Q_s$ a polynomial and $r_k=\Res_{z=z_k}F$, and using
\cref{lem:moment-sums}(a) $2$-adically for $Q_s$, we obtain the identity in $\Q_2[X]$
\begin{multline}\label{eq:H-two-adic}
  H(X)=\sum_{m\ge1}w^m\bigl[\vvec(2m)\ \vvec'(2m)\bigr]
  \begin{pmatrix}F'(2m)&F(2m)\\F(2m)&0\end{pmatrix}
  \bigl[\vvec(2m)\ \vvec'(2m)\bigr]^{\mathsf T}\\
  +(X-\xi_2)\sum_{k<K}\gamma_k\,\vvec(z_k)\vvec(z_k)^{\mathsf T},
\end{multline}
with $\gamma_k=-\sigma^ky^{2k+1}r_k$ and the series converging entrywise in $\Q_2[X]$ for the
Gauss norm. The even sites $2m$ with $1\le m\le A$ contribute nothing, since $F$ and $F'$
vanish there.

The weights have the following valuations. Since $z_k-z_l=2(l-k)$ and $z_k-2j$ is odd,
\begin{equation}\label{eq:gamma-two}
  v_2(\gamma_k)=-(2k+1)-(K-1)-v_2(k!)-v_2((K-1-k)!)\ \ge\ -2k-2K+1 .
\end{equation}
(The inequality holds because $v_2(k!)+v_2((K-1-k)!)\le v_2((K-1)!)\le K-2$, which gives even
$-2k-2K+2$; we keep the weaker form, which is all that is used below.)
For $m>A$ the integer $P_K(2m)$ is odd and
$N^{(2)}(2m)=2^{2A}\bigl(A!\binom{m-1}A\bigr)^2$, so
\[
  v_2(F(2m))\ge1+2A+2v_2(A!)\ge 4A-2\floor{\log_2(A+1)}-1 ,
\]
because $v_2(A!)=A-s_2(A)$ and $s_2(A)\le\floor{\log_2(A+1)}$ (the true bound is
$4A+1-2\floor{\log_2(A+1)}$; we keep the weaker form).
The logarithmic derivative $F'/F=1/z+\sum_j2/(z-2j)-\sum_k1/(z-z_k)$ has valuation at least
$-1-\floor{\log_2m}$ at $z=2m$. Consequently
\begin{equation}\label{eq:nu-m}
  \nu_m:=\min\bigl(v_2F(2m),\,v_2F'(2m)\bigr)\ \ge\ 4A-2\floor{\log_2(A+1)}-2-\floor{\log_2 m}
  \qquad(m>A),
\end{equation}
and $\nu_m\ge0$ when $A=0$, because then $F(2m)$ and $F'(2m)$ are $2$-integral.

\subsection{Minors with nodes of one parity}\label{subsec:parity-minors}

For $t\ge0$ put
\[
  \Pa(t)=\sum_{j<t}\bigl(j+v_2(j!)\bigr)=t(t-1)-\sum_{j<t}s_2(j)=t^2+O(t\log(t+2)).
\]

\begin{lemma}[parity minors]\label{lem:parity-minors}
Let $h\ge2$, $L=\floor{\log_2(h-1)}$, $\varepsilon\in\{0,1\}$, and let $J\subset\{0,\dots,h-1\}$
have $t$ elements. Consider a $t\times t$ matrix whose rows are indexed by $J$ and whose columns
are of the form $(\zeta^i)_{i\in J}$ or $(i\zeta^{i-1})_{i\in J}$ with
$\zeta\in\varepsilon+2\Z$, and let $r$ be the number of columns of the second kind. Then its
determinant has $2$-adic valuation at least $\Pa(t)-r(1+L)$.
\end{lemma}

\begin{proof}
Let $B_j(z)=\prod_{i<j}(z-\varepsilon-2i)$, a monic polynomial in $\Z[z]$ of degree $j$. Then
$z^i=\sum_{j\le i}c_{ij}B_j(z)$ with $c_{ij}\in\Z$ and $c_{ii}=1$. For $\zeta=2s+\varepsilon$,
$s\in\Z$, we have $B_j(\zeta)=2^jj!\binom sj$, which is divisible by $2^{j+v_2(j!)}$. For the
derivative, the identity $\frac{\mathrm d}{\mathrm ds}\binom sj=\sum_{l=1}^{j}
\frac{(-1)^{l-1}}{l}\binom s{j-l}$ (compare coefficients of $T^j$ in
$\partial_s(1+T)^s=\log(1+T)(1+T)^s$) gives $B_j'(\zeta)=2^{j-1}j!\,\partial_s\binom sj$, of
valuation at least $j+v_2(j!)-1-\floor{\log_2j}$. By the Cauchy--Binet formula for the product
of the $t\times h$ matrix $(c_{ij})_{i\in J,\,j<h}$ and the $h\times t$ matrix of the values
$B_j(\zeta)$ or $B'_j(\zeta)$, the determinant is an integral combination of minors with rows
$j_1<\dots<j_t$ of the latter matrix. Since $j\mapsto j+v_2(j!)$ is increasing and
$j_r\ge r-1$, each such minor has valuation at least $\sum_{i<t}(i+v_2(i!))-r(1+L)$.
\end{proof}

\begin{corollary}\label{cor:parity-laplace}
Let $\mathsf E$ be an $h\times h$ matrix whose columns are $P$ columns $\vvec(\zeta)$ with $\zeta$ odd
and $h-P$ columns $\vvec(\zeta)$ or $\vvec'(\zeta)$ with $\zeta$ even, $r$ of them derivatives.
Then $v_2(\det\mathsf E)\ge\Pa(P)+\Pa(h-P)-r(1+L)$.
\end{corollary}

\begin{proof}
Laplace expansion along the odd columns writes $\det\mathsf E$ as a signed sum, over $P$-element sets
$J$ of rows, of products of a $P\times P$ minor with odd nodes and rows $J$ and an
$(h-P)\times(h-P)$ minor with even nodes and the complementary rows. Apply
\cref{lem:parity-minors} to both.
\end{proof}

For nodes of mixed parity the corresponding bound would be the valuation of a Vandermonde
determinant of $t$ consecutive integers, about $t^2/2$; the separation by parity doubles it,
and this is where the constant of \cref{thm:prime-two} comes from.

\subsection{The bound at the prime 2}\label{subsec:two-bound}

\begin{theorem}\label{thm:prime-two}
Let $q$ be odd, $\sigma=\pm1$, and $K\ge h\ge2$, $A\ge0$. Let $\lambda_0\le\lambda_1\le\cdots\le\lambda_{K-1}$ be the exact valuations $v_2(\gamma_k)$, $k<K$, of~\eqref{eq:gamma-two} in increasing order, and $\epsilon_1\le\epsilon_2\le\cdots$ the numbers $2m+\nu_m$, $m>A$, each listed twice, in increasing order. Then
\begin{multline*}
  \vG2(T)\ \ge\ v_2(\nrm)+(K-2A)h-h(h-1)\\
  +\min_{0\le P\le h}\Bigl[2\Pa(P)+2\Pa(h-P)-(h-P)(1+L)+\sum_{i<P}\lambda_i
  +\sum_{i=1}^{h-P}\epsilon_i\Bigr].
\end{multline*} If $K=\kappa h$ and $A=\alpha h$ with $\kappa\ge1$,
$0\le\alpha$ and $d=\kappa-2\alpha>0$ fixed, then
\begin{gather}
  \vG2(T)\ \ge\ h^2\min_{0\le x\le1}f(x)-O(h\log h),\label{eq:two-asymptotic}\\
  f(x)=2(d-1)-4\kappa x+3x^2+\tfrac52(1-x)^2+6\alpha(1-x).\notag
\end{gather}
For $\Gen$ the minimum is $-140/99$, attained at $x=29/33$; for $\Band$ it is $-1847/1100$,
attained at $x=51/55$.
\end{theorem}

\begin{proof}
By~\eqref{eq:Delta-H}, $\vG2(T)=v_2(\nrm)+(K-2A)h-h(h-1)+\vG2(\det H)$. Truncate the sum over
$m$ in~\eqref{eq:H-two-adic} at $m\le M$; the truncated matrices converge to $H$ in the Gauss
norm, so it suffices to bound the determinant of each truncation uniformly in $M$. Write the
truncation as $\mathsf E\mathsf B\mathsf E^{\mathsf T}$, where the columns of $\mathsf E$ are the slots $\vvec(2m)$,
$\vvec'(2m)$ ($A<m\le M$) and $\vvec(z_k)$ ($k<K$), and $\mathsf B$ is block diagonal with the $2\times2$
blocks $w^m\bigl(\begin{smallmatrix}F'&F\\F&0\end{smallmatrix}\bigr)(2m)$ and the $1\times1$
blocks $(X-\xi_2)\gamma_k$. By~\eqref{eq:cauchy-binet}, $\det H$ is a sum of terms indexed by
pairs $(S,S')$ of sets of $h$ slots that select the same poles and the same number $n_m\in\{1,2\}$
of slots at each selected site $2m$.

Fix such a term with $P$ poles and $e=h-P$ even slots. The middle factor has Gauss valuation at
least $\sum_{k\in S}v_2(\gamma_k)+\sum_m n_m(2m+\nu_m)$, since a $1\times1$ minor of the block at
$2m$ is $w^mF'(2m)$ or $w^mF(2m)$ and the $2\times2$ minor is $-w^{2m}F(2m)^2$; and
$\vG2(X-\xi_2)=0$. This is at least $\sum_{i<P}\lambda_i+\sum_{i\le e}\epsilon_i$. The
derivative--derivative entry of each block vanishes, so a site carrying one slot on each side
contributes at most one derivative column to the two outer factors together, and a site carrying
two slots contributes two; the total number of derivative columns on both sides is therefore at
most $e$. By \cref{cor:parity-laplace} the two outer factors together have valuation at least
$2\Pa(P)+2\Pa(e)-e(1+L)$. This proves the first inequality.

For the asymptotic form, $v_2(n!)=n-s_2(n)$ gives $v_2(\nrm)=(K-2A)h-h^2+O(h\log h)$ and
$\Pa(t)=t^2+O(h\log h)$ for $t\le h$. The $P$ smallest of the lower
bounds~\eqref{eq:gamma-two} are those with $k=K-1,\dots,K-P$, whose sum is $-4KP+P^2+2P$. For
the even slots, \eqref{eq:nu-m} and $\floor{\log_2m}\le\log_2(2h)+m/h$ give
$2m+\nu_m\ge(2-1/h)m+4A-O(\log h)$; since each site carries at most two slots and the sites
start at $A+1$, the $e$ smallest values sum to at least
$(2-1/h)\sum_{i=1}^{e}(A+\lceil i/2\rceil)+4Ae-O(h\log h)=6Ae+e^2/2-O(h\log h)$. Collecting
terms with $P=xh$, $e=(1-x)h$ gives~\eqref{eq:two-asymptotic}. Finally
$f'(x)=11x-4\kappa-5-6\alpha$ vanishes at $x=(4\kappa+5+6\alpha)/11$, which is $29/33$ for
$(\kappa,\alpha)=(7/6,0)$ and $51/55$ for $(49/40,1/20)$; substituting gives the two values.
\end{proof}

\begin{remark}\label{rem:two-comparison}
For $(\kappa,\alpha)=(1,0)$, Krisman's determinant, \eqref{eq:two-asymptotic} gives
$-\frac{13}{11}h^2$ at $x=9/11$. Krisman proves the bound $-\frac43h^2$
\cite[Lemma~11.1]{Krisman2026chi2}, shows that $-h^2+h$ is attained by the leading coefficient
\cite[Prop.~11.2]{Krisman2026chi2}, and reports a measured $2$-adic loss of about $1.18\log2$ per $h^2$ at $h=64$, in a
normalisation we cannot identify from his text. In ours, the exact values for $q=83$,
$\sigma=1$, $K=h$, $A=0$ are given in \cref{tab:two-adic-square}: the loss of the un-normalised
determinant $\Delta=\det[\mu_X(\varphi u^{i+j})]$ is $1.1758$ per $h^2$ at $h=64$, and that of
the normalised polynomial $T$ of this paper is between $1.00$ and $1.11$ for $24\le h\le64$,
against $\frac{13}{11}\approx1.18$ for the limit of the bound. These are different quantities
and the finite bound is still far from its limit at these sizes, so we draw no conclusion from the
proximity of $1.1758$ to $1.18$. The terms with $P=h$ are the only ones that contribute to the
leading coefficient in $X$, and $f(1)=1-2\kappa-4\alpha$ equals $-1$ at $(\kappa,\alpha)=(1,0)$,
in accordance with his Proposition~11.2. The minimum
of $f$ is attained at an interior $x$, that is, by coefficients of lower degree in $X$. For
$\Gen$ at $q=83$ and $h=6,12,\dots,72$, the finite bound of \cref{thm:prime-two} lies below the
exact value of $\vG2(T)$ by between $3$ and $56$ (\cref{tab:two-adic}); the ratios
$\vG2(T)/h^2$ are not monotone in $h$ and all exceed $-140/99$. This is consistency evidence
only.
\end{remark}
